\documentclass[12pt,a4paper]{article}

\usepackage[utf8]{inputenc}
\usepackage[T1]{fontenc}
\usepackage[spanish,es-nodecimaldot]{babel}
\usepackage{geometry}
\usepackage{graphicx}
\usepackage{float}
\usepackage{hyperref}
\usepackage{enumitem}
\usepackage{xcolor}
\usepackage{listings}
\usepackage{booktabs}
\usepackage{array}
\usepackage{amsmath}

\geometry{margin=2.5cm}

\hypersetup{
  colorlinks=true,
  linkcolor=blue!60!black,
  urlcolor=blue!60!black,
  pdftitle={Tarea 5 - Instancing y Vertex Animation Textures},
  pdfauthor={Diego Barrera Hernández y Brandon Daniel Martínez Alcántara}
}

\definecolor{codebg}{RGB}{20,20,32}
\definecolor{codeframe}{RGB}{60,60,100}
\definecolor{codecomment}{RGB}{100,160,100}
\definecolor{codekw}{RGB}{130,160,255}
\definecolor{codestr}{RGB}{255,180,100}

\lstdefinestyle{glsl}{
  backgroundcolor=\color{codebg},
  basicstyle=\ttfamily\footnotesize\color{white},
  keywordstyle=\color{codekw}\bfseries,
  commentstyle=\color{codecomment},
  stringstyle=\color{codestr},
  frame=single,
  rulecolor=\color{codeframe},
  breaklines=true,
  tabsize=2,
}

\setlength{\parindent}{0pt}
\setlength{\parskip}{0.8em}

% ──────────────────────────────────────────────
\begin{document}

% ── PORTADA ───────────────────────────────────
\begin{titlepage}
  \centering
  \vspace*{1.5cm}

  {\Large Universidad Autónoma Metropolitana\par}
  \vspace{0.4cm}
  {\large Licenciatura en tecnologias y sistemas de la información.\par}

  \vspace{2.5cm}
  {\Huge\bfseries Seminario de Tecnologías II\par}
  \vspace{0.6cm}
  {\LARGE Tarea 5\par}
  \vspace{0.4cm}
  {\large Instancing y Vertex Animation Textures\par}

  \vspace{3cm}
  {\large Diego Barrera Hernández\par}
  \vspace{0.3cm}
  {\large Brandon Daniel Martínez Alcántara\par}

  \vfill
  {\large Ciudad de México, 2025\par}
\end{titlepage}

% ── ÍNDICE ────────────────────────────────────
\tableofcontents
\newpage

% ─────────────────────────────────────────────
\section{Objetivo}

El objetivo de esta tarea es comparar dos enfoques para renderizar grandes cantidades de
agentes animados en tiempo real: \textbf{Instancing con actualización de matrices en CPU}
(Parte~D) y \textbf{Vertex Animation Textures horneadas} (Parte~E). Ambas partes usan
\texttt{THREE.InstancedMesh} para emitir un único draw call independientemente del número
de instancias, pero difieren radicalmente en dónde vive el costo de animación.

El proyecto se divide en dos partes:

\begin{itemize}[leftmargin=*]
  \item \textbf{Parte D (50 pts):} \texttt{InstancedMesh} con movimiento simulado —
        las matrices de cada instancia se recalculan en CPU cada frame.
  \item \textbf{Parte E (50 pts):} \texttt{InstancedMesh} + Bone VAT —
        la animación de esqueleto se hornea en una \texttt{DataTexture} y el
        vertex shader la reproduce en GPU sin ningún trabajo en CPU por frame.
\end{itemize}

La demo desplegada está disponible en:

\begin{center}
  \url{https://tareas-seminario.vercel.app/}
\end{center}

% ─────────────────────────────────────────────
\section{Escena implementada: Cardumen de peces}

Ambas partes trabajan sobre la misma escena: un cardumen de peces payaso
(\textit{clown fish}) renderizados en un entorno subacuático oscuro con niebla
exponencial. El modelo \texttt{clown\_fish\_low\_poly\_animated.glb} es un
\texttt{SkinnedMesh} con animación de esqueleto (\texttt{swim}) cargado mediante
\texttt{GLTFLoader}.

La escena incluye:
\begin{itemize}[leftmargin=*]
  \item Fondo oscuro (\texttt{\#071521}) con niebla exponencial (\textit{density} = 0.014).
  \item Luz hemisférica (cielo azulado / fondo oscuro) + luz direccional cálida.
  \item Suelo circular con \texttt{MeshStandardMaterial} rugoso.
  \item Cámara con \texttt{OrbitControls} (damping 0.08) posicionada en $(0,\;16,\;48)$.
  \item Límites de movimiento $\pm 30$ en $X$ y $Z$, $[0.6,\;11]$ en $Y$.
\end{itemize}

El modelo se normaliza al cargar: se centra en el origen, se escala uniformemente
para que su mayor dimensión mida 2.6 unidades y se rota 180° en $Y$ para que los
peces apunten en la dirección de avance.

% ─────────────────────────────────────────────
\section{Parte D — Instancing con InstancedMesh}

\subsection{El concepto de Instancing}

\textit{Instancing} permite renderizar $N$ copias de la misma geometría con un único
draw call. En lugar de llamar $N$ veces a \texttt{drawElements}, la GPU recibe la
geometría una vez y un buffer de matrices que describe la posición, rotación y escala
de cada copia. Esto elimina el overhead de CPU de $N$ draw calls y es especialmente
eficiente cuando las instancias comparten la misma malla y material.

En Three.js, \texttt{THREE.InstancedMesh} encapsula este mecanismo: acepta una
\texttt{BufferGeometry}, un \texttt{Material} y el número máximo de instancias, y
expone \texttt{setMatrixAt(i, matrix)} para actualizar la transformación de cada instancia.

\subsection{Inicialización y pre-asignación}

El sistema pre-asigna \textbf{2 000 instancias} al crear el \texttt{InstancedMesh},
lo que permite cambiar la cantidad activa en tiempo de ejecución sin reasignar buffers
de GPU. El estado de simulación de cada instancia vive en \texttt{Float32Array}
paralelos:

\begin{lstlisting}[style=glsl, language={}]
const MAX_INSTANCES = 2000;
const posX    = new Float32Array(MAX_INSTANCES);
const posY    = new Float32Array(MAX_INSTANCES);
const posZ    = new Float32Array(MAX_INSTANCES);
const heading = new Float32Array(MAX_INSTANCES);
const speed   = new Float32Array(MAX_INSTANCES);
const phase   = new Float32Array(MAX_INSTANCES);
const scaleArr = new Float32Array(MAX_INSTANCES);
const turnRate = new Float32Array(MAX_INSTANCES);
\end{lstlisting}

Los arreglos de tipo (\texttt{Float32Array}) se eligen deliberadamente en lugar de
arrays JavaScript normales: son datos contiguos en memoria y permiten iteraciones
más rápidas en el bucle de actualización al evitar la indiferencia de tipo del motor JS.

La carga del modelo elimina los atributos \texttt{skinIndex} y \texttt{skinWeight}
(innecesarios en esta parte) y crea un único \texttt{InstancedMesh}:

\begin{lstlisting}[style=glsl, language={}]
geometry.deleteAttribute('skinIndex');
geometry.deleteAttribute('skinWeight');
instancedMesh = new THREE.InstancedMesh(geometry, material, MAX_INSTANCES);
instancedMesh.instanceMatrix.setUsage(THREE.DynamicDrawUsage);
instancedMesh.count = activeCount;  // comienza en 500
\end{lstlisting}

\texttt{DynamicDrawUsage} indica a la GPU que el buffer de matrices se actualizará
frecuentemente, lo que puede mejorar el rendimiento en algunas implementaciones WebGL.

\subsection{Bucle de actualización de instancias}

Cada frame, la función \texttt{updateInstances(dt, elapsed)} itera sobre las
\texttt{activeCount} instancias y calcula la nueva posición y orientación:

\begin{lstlisting}[style=glsl, language={}]
for (let i = 0; i < count; i++) {
  // Giro suave: componente determinista + oscilacion senoidal
  heading[i] += turnRate[i] * dt * 0.2
             +  Math.sin(elapsed * 0.5 + phase[i]) * dt * 0.3;

  // Avance en la direccion actual
  posX[i] += Math.cos(heading[i]) * speed[i] * dt;
  posZ[i] += Math.sin(heading[i]) * speed[i] * dt;
  posY[i] += Math.sin(elapsed * 1.3 + phase[i]) * dt * 0.6;

  // Wrapping toroidal en los limites del volumen
  if (posX[i] >  BOUNDS.x) posX[i] = -BOUNDS.x;
  if (posX[i] < -BOUNDS.x) posX[i] =  BOUNDS.x;
  // ... idem Z, Y con clamp

  dummy.position.set(posX[i], posY[i], posZ[i]);
  dummy.rotation.set(Math.sin(elapsed*4 + phase[i])*0.08,
                     heading[i], 0);
  dummy.scale.setScalar(scaleArr[i]);
  dummy.updateMatrix();
  instancedMesh.setMatrixAt(i, dummy.matrix);
}
instancedMesh.instanceMatrix.needsUpdate = true;
\end{lstlisting}

Al final del bucle, \texttt{instanceMatrix.needsUpdate = true} sube el buffer de
matrices completo a la GPU en el siguiente frame. Esto es el costo dominante de la
Parte~D: $O(N)$ operaciones en CPU por frame y una transferencia de $N \times 64$
bytes (cada \texttt{mat4}) al bus de la GPU.

\subsection{Simulación de natación}

El movimiento visible de los peces combina tres componentes, todas calculadas en CPU:

\begin{center}
\begin{tabular}{ll}
  \toprule
  Componente & Mecanismo \\
  \midrule
  Avance horizontal & \texttt{cos/sin(heading)} $\times$ \texttt{speed} $\times$ \texttt{dt} \\
  Oscilación vertical & $\sin(elapsed \times 1.3 + phase_i) \times 0.6 \times dt$ \\
  Bamboleo de cabeza & Rotación en $X$: $\sin(elapsed \times 4 + phase_i) \times 0.08$ \\
  Giro & \texttt{turnRate}$_i$ + perturbación senoidal lenta \\
  \bottomrule
\end{tabular}
\end{center}

La propiedad \texttt{phase}$_i$ inicializada aleatoriamente garantiza que cada pez
esté en un punto distinto de su ciclo de movimiento desde el primer frame.

\subsection{Control de cantidad}

Un slider en el panel lateral cambia \texttt{instancedMesh.count} sin reconstruir
nada: Three.js simplemente ignora las instancias con índice mayor al \texttt{count}
activo. Esto permite observar en tiempo real cómo el FPS cae a medida que aumenta
el número de instancias activas, ilustrando la complejidad $O(N)$ del enfoque.

% ─────────────────────────────────────────────
\section{Parte E — Bone VAT (Vertex Animation Texture)}

\subsection{El concepto de VAT}

Una \textit{Vertex Animation Texture} (VAT) o, en este caso, \textit{Bone VAT},
consiste en hornear los resultados de una animación de esqueleto a una textura y
reproducirla en el vertex shader. En lugar de ejecutar el \texttt{AnimationMixer}
de Three.js en CPU cada frame (que calcula e interpola la pose de todos los huesos),
el shader simplemente \textbf{muestrea la textura} para obtener la matriz de cada hueso
en el instante deseado.

El resultado es que el costo de animación pasa de $O(N \times B)$ en CPU
($N$ instancias, $B$ huesos) a un costo fijo de GPU que depende solo de la
complejidad del shader, no del número de instancias.

\subsection{Horneado: \texttt{bakeBoneVAT}}

La función \texttt{bakeBoneVAT(gltf, clipName)} recorre el clip de animación a
\textbf{24 fps} y almacena la pose de cada hueso en cada frame:

\begin{enumerate}[leftmargin=*]
  \item Se localiza el \texttt{SkinnedMesh} y su \texttt{Skeleton} dentro de la escena GLTF.
  \item Se crea un \texttt{AnimationMixer} efímero y se hace avanzar al clip frame a frame
        con \texttt{mixer.setTime(t)} + \texttt{skeleton.update()}.
  \item En cada frame $f$, \texttt{skeleton.boneMatrices} contiene el arreglo de matrices
        $4\times4$ de todos los huesos ya multiplicadas por la bindMatrixInverse
        (listas para el skinning). Estas 16 floats por hueso se copian a un
        \texttt{Float32Array} de tamaño \texttt{numFrames × numBones × 16}.
  \item Los datos se suben como \texttt{DataTexture} de tipo \texttt{FloatType},
        con \texttt{NearestFilter} (para evitar interpolación no deseada entre huesos)
        y sin mipmaps.
\end{enumerate}

\begin{lstlisting}[style=glsl, language={}]
const BAKE_FPS  = 24;
const numFrames = Math.round(clip.duration * BAKE_FPS);
const data = new Float32Array(numFrames * numBones * 16);

for (let f = 0; f < numFrames; f++) {
  mixer.setTime((f / (numFrames - 1)) * clip.duration);
  gltf.scene.updateMatrixWorld(true);
  skeleton.update();
  data.set(skeleton.boneMatrices, f * numBones * 16);
}

const texture = new THREE.DataTexture(
  data, numBones * 4, numFrames,
  THREE.RGBAFormat, THREE.FloatType
);
texture.magFilter = THREE.NearestFilter;
texture.needsUpdate = true;
\end{lstlisting}

La textura tiene \textbf{ancho} = $4 \times \texttt{numBones}$ texeles
(cada mat4 ocupa 4 texeles RGBA de 4 floats cada uno) y
\textbf{alto} = \texttt{numFrames} filas. Una vez creada, el mixer se descarta:
ya no se necesita en ningún frame posterior.

\subsection{Offset de tiempo por instancia}

Para que las 500 instancias naden de forma desincronizada sin ningún costo adicional
por frame, se añade un atributo \texttt{timeOffset} como
\texttt{InstancedBufferAttribute}:

\begin{lstlisting}[style=glsl, language={}]
const timeOffsets = new Float32Array(INSTANCE_COUNT);
for (let i = 0; i < INSTANCE_COUNT; i++) {
  timeOffsets[i] = Math.random() * baked.duration;
}
geometry.setAttribute('timeOffset',
  new THREE.InstancedBufferAttribute(timeOffsets, 1));
\end{lstlisting}

Este buffer se sube a la GPU una única vez al cargar. El vertex shader usa
\texttt{timeOffset} (por instancia) más el uniforme \texttt{time} (global)
para calcular el frame de animación de cada instancia independientemente:

\begin{lstlisting}[style=glsl, language={}]
float localTime = mod(time + timeOffset, duration);
float framef    = (localTime / duration) * numFrames;
float f0        = floor(framef);
float f1        = mod(f0 + 1.0, numFrames);
float frac      = framef - f0;
\end{lstlisting}

\subsection{Vertex shader — lectura de la textura VAT}

La función auxiliar \texttt{getBoneMatrix} reconstruye una mat4 a partir de cuatro
texeles consecutivos de la fila correspondiente al frame:

\begin{lstlisting}[style=glsl, language={}]
mat4 getBoneMatrix(float boneIndex, float frame) {
  float texWidth = numBones * 4.0;
  float u0 = (boneIndex * 4.0 + 0.5) / texWidth;
  float du = 1.0 / texWidth;
  float v  = (frame + 0.5) / numFrames;
  vec4 c0 = texture2D(boneTexture, vec2(u0,           v));
  vec4 c1 = texture2D(boneTexture, vec2(u0 + du,      v));
  vec4 c2 = texture2D(boneTexture, vec2(u0 + 2.0*du,  v));
  vec4 c3 = texture2D(boneTexture, vec2(u0 + 3.0*du,  v));
  return mat4(c0, c1, c2, c3);
}
\end{lstlisting}

Con las matrices de los frames $f_0$ y $f_1$ se calcula el skinning en ambos
y se interpola linealmente con \texttt{mix} para conseguir una animación suave
entre fotogramas horneados:

\begin{lstlisting}[style=glsl, language={}]
vec4 skinVertex = bindMatrix * vec4(position, 1.0);

mat4 bx0 = getBoneMatrix(skinIndex.x, f0);
// ... (by0, bz0, bw0 idem para los otros 3 huesos)
mat4 bx1 = getBoneMatrix(skinIndex.x, f1);
// ... (by1, bz1, bw1)

vec4 pos0 = bx0*skinVertex*skinWeight.x + by0*skinVertex*skinWeight.y
          + bz0*skinVertex*skinWeight.z + bw0*skinVertex*skinWeight.w;
vec4 pos1 = bx1*skinVertex*skinWeight.x + by1*skinVertex*skinWeight.y
          + bz1*skinVertex*skinWeight.z + bw1*skinVertex*skinWeight.w;

vec4 skinnedPos = mix(pos0, pos1, frac);
vec3 localPos   = (bindMatrixInverse * skinnedPos).xyz;
\end{lstlisting}

La posición final en clip space combina la transformación de instancia
(\texttt{instanceMatrix}, fijada una sola vez al cargar) con la posición de modelo:

\begin{lstlisting}[style=glsl, language={}]
vec4 worldPos = modelMatrix * instanceMatrix * vec4(localPos, 1.0);
gl_Position   = projectionMatrix * viewMatrix * worldPos;
\end{lstlisting}

\subsection{Fragment shader — iluminación Blinn-Phong con textura}

El fragment shader aplica la textura de color del modelo con iluminación
Blinn-Phong usando el half-vector $\mathbf{H} = \text{normalize}(\mathbf{L} + \mathbf{V})$:

\begin{lstlisting}[style=glsl, language={}]
vec3 N = normalize(vNormal);
vec3 L = normalize(-lightDir);
vec3 V = normalize(viewPos - vWorldPos);
vec3 H = normalize(L + V);

vec4 base = texture2D(map, vUv);
float diff = max(dot(N, L), 0.0);
float spec = pow(max(dot(N, H), 0.0), 32.0);

vec3 color = base.rgb * 0.35            // termino ambiental
           + base.rgb * diff * lightColor  // difuso
           + vec3(0.18) * spec;            // especular

gl_FragColor = vec4(color, base.a);
\end{lstlisting}

\subsection{Trabajo de CPU por frame}

Una vez terminada la carga, el único trabajo de CPU por frame es:

\begin{lstlisting}[style=glsl, language={}]
if (material) material.uniforms.time.value = elapsed;
\end{lstlisting}

Una asignación de un flotante. Las matrices de instancia \textbf{nunca se tocan}
tras la inicialización: los peces no se mueven por el espacio en la Parte~E
(cada uno permanece en su posición fija), pero su ciclo de natación es visible
porque el vertex shader los deforma desde la textura VAT.

% ─────────────────────────────────────────────
\section{Estructura del proyecto}

\begin{verbatim}
Tarea_3/
+-- partD/
|   +-- index.html          (interfaz: slider de cantidad, stats)
|   +-- main.js             (punto de entrada)
+-- partE/
|   +-- index.html          (interfaz: stats del VAT)
|   +-- main.js             (punto de entrada)
+-- src/
    +-- shared/
    |   +-- FpsMeter.js     (contador de frames en pantalla)
    +-- instancing/
    |   +-- FishSwarm.js    (Parte D: InstancedMesh + sim CPU)
    +-- vat/
        +-- BoneVAT.js      (horneado de clip a DataTexture)
        +-- FishSwarmVAT.js (Parte E: InstancedMesh + ShaderMaterial)
        +-- boneVatShader.js (vertex y fragment GLSL)
\end{verbatim}

% ─────────────────────────────────────────────
\section{Comparativa de los dos enfoques}

\begin{center}
\begin{tabular}{lcc}
  \toprule
  Aspecto & Parte D (Instancing) & Parte E (Bone VAT) \\
  \midrule
  Draw calls por frame    & 1                       & 1 \\
  Costo de animación      & $O(N)$ en CPU           & $O(1)$ en CPU \\
  Trabajo CPU por frame   & $N$ matrices actualizadas & 1 float (\texttt{time}) \\
  Transferencia GPU/frame & $N \times 64$ bytes      & 4 bytes (uniforme) \\
  Desincronización        & \texttt{phase}$_i$ + sim & \texttt{timeOffset}$_i$ \\
  Skinning                & Ninguno (forma estática) & GPU, vertex shader \\
  Variación de cantidad   & Slider 100--2000        & Fijo en 500 \\
  Modelo usado            & Solo la malla (sin rig) & Malla + rig + clip \\
  \bottomrule
\end{tabular}
\end{center}

La diferencia clave es que en la Parte~D el movimiento de cada pez se simula
matemáticamente en CPU sin skinning real, mientras que la Parte~E reproduce
la animación de esqueleto original del modelo sin ningún trabajo de CPU por frame.

% ─────────────────────────────────────────────
\section{Despliegue}

El proyecto se construye con \textbf{Vite} (\texttt{npm run build}) generando un
sitio estático en \texttt{dist/}. Está desplegado en \textbf{Vercel} con
configuración automática detectada por el preset de Vite:

\begin{itemize}[leftmargin=*]
  \item \textbf{Build command:} \texttt{npm run build}
  \item \textbf{Output directory:} \texttt{dist}
  \item \textbf{Root directory:} \texttt{Tarea\_3}
\end{itemize}

URL pública: \url{https://tareas-seminario.vercel.app/}

Las Partes~D y~E usan \texttt{THREE.WebGLRenderer} estándar, por lo que funcionan
en cualquier navegador moderno con soporte WebGL2 (Chrome 56+, Firefox 51+, Safari 15+).

% ─────────────────────────────────────────────
\section{Conclusiones}

Implementar el mismo cardumen de peces con dos estrategias de animación distintas
permite cuantificar con precisión el trade-off entre flexibilidad y rendimiento:

\begin{itemize}[leftmargin=*]
  \item \textbf{Instancing puro (Parte~D)} es sencillo de implementar: basta con
        mantener arreglos de estado por instancia y actualizar las matrices cada frame.
        Su limitación es que el costo escala linealmente con el número de instancias
        activas, lo que pone un techo práctico de rendimiento visible en tiempo real
        con el slider.

  \item \textbf{Bone VAT (Parte~E)} requiere un paso de preprocesamiento (el horneado)
        y un shader personalizado más complejo, pero la recompensa es que el costo de
        animación en CPU por frame es constante e independiente de cuántas instancias
        haya. Esto hace al enfoque VAT ideal para efectos de masas donde se necesitan
        cientos o miles de personajes animados simultáneamente.

  \item Ambas partes mantienen \textbf{un único draw call} gracias a
        \texttt{InstancedMesh}, confirmando que el cuello de botella en la Parte~D
        está en la CPU (actualización de matrices), no en la GPU.

  \item El atributo \texttt{timeOffset} de la Parte~E demuestra que la desincronización
        de instancias es trivial con VAT: el offset se sube una vez y la GPU se encarga
        del resto sin ningún dato adicional por frame.
\end{itemize}

\end{document}
